\documentclass[a4paper]{article}
\usepackage{tikz}
\usetikzlibrary{shapes.geometric}
\usetikzlibrary{calc}
\usepackage{amsmath}
\usepackage{mathtools}
\usepackage{amsthm}
\usepackage{multicol}

\usepackage{fullpage}
\usepackage{enumitem}
\setlist[itemize]{noitemsep}

\usepackage{natbib}
\setlength{\bibsep}{0.0pt}

\newcommand{\cac}{olive}
\newcommand{\cbc}{teal}
\newcommand{\ccc}{orange}
\newcommand{\cdc}{purple}

\usepackage{mdwlist}
\renewenvironment{description}%
{
\begin{basedescript}{
\desclabelstyle{\nextlinelabel}
\renewcommand{\makelabel}[1]{%
\parbox[b]{\columnwidth}{\bfseries##1}%
}%
\desclabelwidth{2em}}}
{
\end{basedescript}
}

\usepackage{titlesec}
\titleformat*{\section}{\fontsize{10pt}{10pt}\bfseries\selectfont}

\newtheorem{theorem}{Theorem}



\begin{document}



\begin{center}
    
{ \bf \large
The chromatic number  of the associahedron: simple and logarithmic
}

\vspace{5pt}
Benjamin Aram Berendsohn \\
{\it Max Planck Institute for Informatics}

\vspace{5pt}
Jean Cardinal and
John Iacono \\ {\it Université libre de Bruxelles}

\vspace{5pt}
László Kozma \\{\it Technische Universität Dresden}\\%

\end{center}




\begin{abstract}
  Addario-Berry, Reed,  Scott, and Wood  (JoCG 2026) proved that the $n$-dimensional associahedron has chromatic number at most $22500\cdot\log_3 n+O(1)$. We present a much simpler proof that the chromatic number of the $n-1$-dimensional associahedron is at most $6 \log_2 n + 18$. Since this paper's appearance as a preliminary abstract in the Japan Conference on Discrete and Computational Geometry, Graphs, and Games held from September 7-10, 2026, the bound has been improved to $O(\log \log n)$ by Oum and Wood (arXiv September 30, 2026). It is hoped that the simple method here could give insights to further improvements.
\end{abstract}

\begin{figure}[h]
\begin{center}
\resizebox{10cm}{!}{% 42 binary search trees on 1..5, 84 rotations
\definecolor{assocA}{RGB}{204,224,246}
\definecolor{assocB}{RGB}{208,237,213}
\definecolor{assocC}{RGB}{250,240,197}
\definecolor{assocD}{RGB}{250,214,218}
\definecolor{assocE}{RGB}{226,215,243}
\begin{tikzpicture}[
  te/.style={line width=0.35pt,black!75},
  bn/.style={circle,draw,line width=0.3pt,fill=white,inner sep=0pt},
  gedge/.style={line width=0.50pt,black!50},
  rlab/.style={circle,fill=black,text=white,inner sep=0pt,outer sep=0pt,
      minimum size=3.60mm,font=\tiny\sffamily\bfseries},
  vbox/.style={rectangle,rounded corners=3.40pt,draw=black!35,
      line width=0.3pt,inner sep=0.050cm,outer sep=0pt}]
\pgfdeclarelayer{bg}\pgfsetlayers{bg,main}
\node[vbox,fill=assocA] (v0) at (4.2565,4.3527) {\begin{tikzpicture}[x=0.340cm,y=0.300cm]\path (0.50,0.55) rectangle (5.50,-4.55);\draw[te] (1,0)--(2,-1);\draw[te] (2,-1)--(3,-2);\draw[te] (3,-2)--(4,-3);\draw[te] (4,-3)--(5,-4);\node[bn,minimum size=2.70mm,font=\tiny] at (1,0) {1};\node[bn,minimum size=2.70mm,font=\tiny] at (2,-1) {2};\node[bn,minimum size=2.70mm,font=\tiny] at (3,-2) {3};\node[bn,minimum size=2.70mm,font=\tiny] at (4,-3) {4};\node[bn,minimum size=2.70mm,font=\tiny] at (5,-4) {5};\end{tikzpicture}};
\node[vbox,fill=assocC] (v1) at (5.9735,6.6004) {\begin{tikzpicture}[x=0.340cm,y=0.300cm]\path (0.50,0.55) rectangle (5.50,-4.55);\draw[te] (1,0)--(2,-1);\draw[te] (2,-1)--(3,-2);\draw[te] (3,-2)--(5,-3);\draw[te] (5,-3)--(4,-4);\node[bn,minimum size=2.70mm,font=\tiny] at (1,0) {1};\node[bn,minimum size=2.70mm,font=\tiny] at (2,-1) {2};\node[bn,minimum size=2.70mm,font=\tiny] at (3,-2) {3};\node[bn,minimum size=2.70mm,font=\tiny] at (5,-3) {5};\node[bn,minimum size=2.70mm,font=\tiny] at (4,-4) {4};\end{tikzpicture}};
\node[vbox,fill=assocD] (v2) at (5.1418,1.5118) {\begin{tikzpicture}[x=0.340cm,y=0.300cm]\path (0.50,0.55) rectangle (5.50,-3.55);\draw[te] (1,0)--(2,-1);\draw[te] (2,-1)--(4,-2);\draw[te] (4,-2)--(3,-3);\draw[te] (4,-2)--(5,-3);\node[bn,minimum size=2.70mm,font=\tiny] at (1,0) {1};\node[bn,minimum size=2.70mm,font=\tiny] at (2,-1) {2};\node[bn,minimum size=2.70mm,font=\tiny] at (4,-2) {4};\node[bn,minimum size=2.70mm,font=\tiny] at (3,-3) {3};\node[bn,minimum size=2.70mm,font=\tiny] at (5,-3) {5};\end{tikzpicture}};
\node[vbox,fill=assocA] (v3) at (11.4304,0.0713) {\begin{tikzpicture}[x=0.340cm,y=0.300cm]\path (0.50,0.55) rectangle (5.50,-4.55);\draw[te] (1,0)--(2,-1);\draw[te] (2,-1)--(5,-2);\draw[te] (5,-2)--(3,-3);\draw[te] (3,-3)--(4,-4);\node[bn,minimum size=2.70mm,font=\tiny] at (1,0) {1};\node[bn,minimum size=2.70mm,font=\tiny] at (2,-1) {2};\node[bn,minimum size=2.70mm,font=\tiny] at (5,-2) {5};\node[bn,minimum size=2.70mm,font=\tiny] at (3,-3) {3};\node[bn,minimum size=2.70mm,font=\tiny] at (4,-4) {4};\end{tikzpicture}};
\node[vbox,fill=assocC] (v4) at (8.6555,-1.0052) {\begin{tikzpicture}[x=0.340cm,y=0.300cm]\path (0.50,0.55) rectangle (5.50,-4.55);\draw[te] (1,0)--(2,-1);\draw[te] (2,-1)--(5,-2);\draw[te] (5,-2)--(4,-3);\draw[te] (4,-3)--(3,-4);\node[bn,minimum size=2.70mm,font=\tiny] at (1,0) {1};\node[bn,minimum size=2.70mm,font=\tiny] at (2,-1) {2};\node[bn,minimum size=2.70mm,font=\tiny] at (5,-2) {5};\node[bn,minimum size=2.70mm,font=\tiny] at (4,-3) {4};\node[bn,minimum size=2.70mm,font=\tiny] at (3,-4) {3};\end{tikzpicture}};
\node[vbox,fill=assocC] (v5) at (-0.0506,3.5508) {\begin{tikzpicture}[x=0.340cm,y=0.300cm]\path (0.50,0.55) rectangle (5.50,-3.55);\draw[te] (1,0)--(3,-1);\draw[te] (3,-1)--(2,-2);\draw[te] (3,-1)--(4,-2);\draw[te] (4,-2)--(5,-3);\node[bn,minimum size=2.70mm,font=\tiny] at (1,0) {1};\node[bn,minimum size=2.70mm,font=\tiny] at (3,-1) {3};\node[bn,minimum size=2.70mm,font=\tiny] at (2,-2) {2};\node[bn,minimum size=2.70mm,font=\tiny] at (4,-2) {4};\node[bn,minimum size=2.70mm,font=\tiny] at (5,-3) {5};\end{tikzpicture}};
\node[vbox,fill=assocD] (v6) at (1.8149,5.4227) {\begin{tikzpicture}[x=0.340cm,y=0.300cm]\path (0.50,0.55) rectangle (5.50,-3.55);\draw[te] (1,0)--(3,-1);\draw[te] (3,-1)--(2,-2);\draw[te] (3,-1)--(5,-2);\draw[te] (5,-2)--(4,-3);\node[bn,minimum size=2.70mm,font=\tiny] at (1,0) {1};\node[bn,minimum size=2.70mm,font=\tiny] at (3,-1) {3};\node[bn,minimum size=2.70mm,font=\tiny] at (2,-2) {2};\node[bn,minimum size=2.70mm,font=\tiny] at (5,-2) {5};\node[bn,minimum size=2.70mm,font=\tiny] at (4,-3) {4};\end{tikzpicture}};
\node[vbox,fill=assocC] (v7) at (-2.5961,-2.9439) {\begin{tikzpicture}[x=0.340cm,y=0.300cm]\path (0.50,0.55) rectangle (5.50,-3.55);\draw[te] (1,0)--(4,-1);\draw[te] (4,-1)--(2,-2);\draw[te] (2,-2)--(3,-3);\draw[te] (4,-1)--(5,-2);\node[bn,minimum size=2.70mm,font=\tiny] at (1,0) {1};\node[bn,minimum size=2.70mm,font=\tiny] at (4,-1) {4};\node[bn,minimum size=2.70mm,font=\tiny] at (2,-2) {2};\node[bn,minimum size=2.70mm,font=\tiny] at (3,-3) {3};\node[bn,minimum size=2.70mm,font=\tiny] at (5,-2) {5};\end{tikzpicture}};
\node[vbox,fill=assocD] (v8) at (-5.0890,-3.8209) {\begin{tikzpicture}[x=0.340cm,y=0.300cm]\path (0.50,0.55) rectangle (5.50,-3.55);\draw[te] (1,0)--(4,-1);\draw[te] (4,-1)--(3,-2);\draw[te] (3,-2)--(2,-3);\draw[te] (4,-1)--(5,-2);\node[bn,minimum size=2.70mm,font=\tiny] at (1,0) {1};\node[bn,minimum size=2.70mm,font=\tiny] at (4,-1) {4};\node[bn,minimum size=2.70mm,font=\tiny] at (3,-2) {3};\node[bn,minimum size=2.70mm,font=\tiny] at (2,-3) {2};\node[bn,minimum size=2.70mm,font=\tiny] at (5,-2) {5};\end{tikzpicture}};
\node[vbox,fill=assocD] (v9) at (6.2162,-5.7010) {\begin{tikzpicture}[x=0.340cm,y=0.300cm]\path (0.50,0.55) rectangle (5.50,-4.55);\draw[te] (1,0)--(5,-1);\draw[te] (5,-1)--(2,-2);\draw[te] (2,-2)--(3,-3);\draw[te] (3,-3)--(4,-4);\node[bn,minimum size=2.70mm,font=\tiny] at (1,0) {1};\node[bn,minimum size=2.70mm,font=\tiny] at (5,-1) {5};\node[bn,minimum size=2.70mm,font=\tiny] at (2,-2) {2};\node[bn,minimum size=2.70mm,font=\tiny] at (3,-3) {3};\node[bn,minimum size=2.70mm,font=\tiny] at (4,-4) {4};\end{tikzpicture}};
\node[vbox,fill=assocA] (v10) at (6.9531,-9.3425) {\begin{tikzpicture}[x=0.340cm,y=0.300cm]\path (0.50,0.55) rectangle (5.50,-4.55);\draw[te] (1,0)--(5,-1);\draw[te] (5,-1)--(2,-2);\draw[te] (2,-2)--(4,-3);\draw[te] (4,-3)--(3,-4);\node[bn,minimum size=2.70mm,font=\tiny] at (1,0) {1};\node[bn,minimum size=2.70mm,font=\tiny] at (5,-1) {5};\node[bn,minimum size=2.70mm,font=\tiny] at (2,-2) {2};\node[bn,minimum size=2.70mm,font=\tiny] at (4,-3) {4};\node[bn,minimum size=2.70mm,font=\tiny] at (3,-4) {3};\end{tikzpicture}};
\node[vbox,fill=assocC] (v11) at (0.4739,-3.4052) {\begin{tikzpicture}[x=0.340cm,y=0.300cm]\path (0.50,0.55) rectangle (5.50,-3.55);\draw[te] (1,0)--(5,-1);\draw[te] (5,-1)--(3,-2);\draw[te] (3,-2)--(2,-3);\draw[te] (3,-2)--(4,-3);\node[bn,minimum size=2.70mm,font=\tiny] at (1,0) {1};\node[bn,minimum size=2.70mm,font=\tiny] at (5,-1) {5};\node[bn,minimum size=2.70mm,font=\tiny] at (3,-2) {3};\node[bn,minimum size=2.70mm,font=\tiny] at (2,-3) {2};\node[bn,minimum size=2.70mm,font=\tiny] at (4,-3) {4};\end{tikzpicture}};
\node[vbox,fill=assocD] (v12) at (-0.8109,-8.8648) {\begin{tikzpicture}[x=0.340cm,y=0.300cm]\path (0.50,0.55) rectangle (5.50,-4.55);\draw[te] (1,0)--(5,-1);\draw[te] (5,-1)--(4,-2);\draw[te] (4,-2)--(2,-3);\draw[te] (2,-3)--(3,-4);\node[bn,minimum size=2.70mm,font=\tiny] at (1,0) {1};\node[bn,minimum size=2.70mm,font=\tiny] at (5,-1) {5};\node[bn,minimum size=2.70mm,font=\tiny] at (4,-2) {4};\node[bn,minimum size=2.70mm,font=\tiny] at (2,-3) {2};\node[bn,minimum size=2.70mm,font=\tiny] at (3,-4) {3};\end{tikzpicture}};
\node[vbox,fill=assocA] (v13) at (-3.0587,-5.9964) {\begin{tikzpicture}[x=0.340cm,y=0.300cm]\path (0.50,0.55) rectangle (5.50,-4.55);\draw[te] (1,0)--(5,-1);\draw[te] (5,-1)--(4,-2);\draw[te] (4,-2)--(3,-3);\draw[te] (3,-3)--(2,-4);\node[bn,minimum size=2.70mm,font=\tiny] at (1,0) {1};\node[bn,minimum size=2.70mm,font=\tiny] at (5,-1) {5};\node[bn,minimum size=2.70mm,font=\tiny] at (4,-2) {4};\node[bn,minimum size=2.70mm,font=\tiny] at (3,-3) {3};\node[bn,minimum size=2.70mm,font=\tiny] at (2,-4) {2};\end{tikzpicture}};
\node[vbox,fill=assocD] (v14) at (3.4758,7.9123) {\begin{tikzpicture}[x=0.340cm,y=0.300cm]\path (0.50,0.55) rectangle (5.50,-3.55);\draw[te] (2,0)--(1,-1);\draw[te] (2,0)--(3,-1);\draw[te] (3,-1)--(4,-2);\draw[te] (4,-2)--(5,-3);\node[bn,minimum size=2.70mm,font=\tiny] at (2,0) {2};\node[bn,minimum size=2.70mm,font=\tiny] at (1,-1) {1};\node[bn,minimum size=2.70mm,font=\tiny] at (3,-1) {3};\node[bn,minimum size=2.70mm,font=\tiny] at (4,-2) {4};\node[bn,minimum size=2.70mm,font=\tiny] at (5,-3) {5};\end{tikzpicture}};
\node[vbox,fill=assocA] (v15) at (6.8620,9.4411) {\begin{tikzpicture}[x=0.340cm,y=0.300cm]\path (0.50,0.55) rectangle (5.50,-3.55);\draw[te] (2,0)--(1,-1);\draw[te] (2,0)--(3,-1);\draw[te] (3,-1)--(5,-2);\draw[te] (5,-2)--(4,-3);\node[bn,minimum size=2.70mm,font=\tiny] at (2,0) {2};\node[bn,minimum size=2.70mm,font=\tiny] at (1,-1) {1};\node[bn,minimum size=2.70mm,font=\tiny] at (3,-1) {3};\node[bn,minimum size=2.70mm,font=\tiny] at (5,-2) {5};\node[bn,minimum size=2.70mm,font=\tiny] at (4,-3) {4};\end{tikzpicture}};
\node[vbox,fill=assocC] (v16) at (2.5153,1.8032) {\begin{tikzpicture}[x=0.340cm,y=0.300cm]\path (0.50,0.55) rectangle (5.50,-2.55);\draw[te] (2,0)--(1,-1);\draw[te] (2,0)--(4,-1);\draw[te] (4,-1)--(3,-2);\draw[te] (4,-1)--(5,-2);\node[bn,minimum size=2.70mm,font=\tiny] at (2,0) {2};\node[bn,minimum size=2.70mm,font=\tiny] at (1,-1) {1};\node[bn,minimum size=2.70mm,font=\tiny] at (4,-1) {4};\node[bn,minimum size=2.70mm,font=\tiny] at (3,-2) {3};\node[bn,minimum size=2.70mm,font=\tiny] at (5,-2) {5};\end{tikzpicture}};
\node[vbox,fill=assocD] (v17) at (8.1239,1.7655) {\begin{tikzpicture}[x=0.340cm,y=0.300cm]\path (0.50,0.55) rectangle (5.50,-3.55);\draw[te] (2,0)--(1,-1);\draw[te] (2,0)--(5,-1);\draw[te] (5,-1)--(3,-2);\draw[te] (3,-2)--(4,-3);\node[bn,minimum size=2.70mm,font=\tiny] at (2,0) {2};\node[bn,minimum size=2.70mm,font=\tiny] at (1,-1) {1};\node[bn,minimum size=2.70mm,font=\tiny] at (5,-1) {5};\node[bn,minimum size=2.70mm,font=\tiny] at (3,-2) {3};\node[bn,minimum size=2.70mm,font=\tiny] at (4,-3) {4};\end{tikzpicture}};
\node[vbox,fill=assocA] (v18) at (5.8276,-1.0643) {\begin{tikzpicture}[x=0.340cm,y=0.300cm]\path (0.50,0.55) rectangle (5.50,-3.55);\draw[te] (2,0)--(1,-1);\draw[te] (2,0)--(5,-1);\draw[te] (5,-1)--(4,-2);\draw[te] (4,-2)--(3,-3);\node[bn,minimum size=2.70mm,font=\tiny] at (2,0) {2};\node[bn,minimum size=2.70mm,font=\tiny] at (1,-1) {1};\node[bn,minimum size=2.70mm,font=\tiny] at (5,-1) {5};\node[bn,minimum size=2.70mm,font=\tiny] at (4,-2) {4};\node[bn,minimum size=2.70mm,font=\tiny] at (3,-3) {3};\end{tikzpicture}};
\node[vbox,fill=assocD] (v19) at (-4.2280,8.1108) {\begin{tikzpicture}[x=0.340cm,y=0.300cm]\path (0.50,0.55) rectangle (5.50,-2.55);\draw[te] (3,0)--(1,-1);\draw[te] (1,-1)--(2,-2);\draw[te] (3,0)--(4,-1);\draw[te] (4,-1)--(5,-2);\node[bn,minimum size=2.70mm,font=\tiny] at (3,0) {3};\node[bn,minimum size=2.70mm,font=\tiny] at (1,-1) {1};\node[bn,minimum size=2.70mm,font=\tiny] at (2,-2) {2};\node[bn,minimum size=2.70mm,font=\tiny] at (4,-1) {4};\node[bn,minimum size=2.70mm,font=\tiny] at (5,-2) {5};\end{tikzpicture}};
\node[vbox,fill=assocA] (v20) at (-0.9582,6.5017) {\begin{tikzpicture}[x=0.340cm,y=0.300cm]\path (0.50,0.55) rectangle (5.50,-2.55);\draw[te] (3,0)--(1,-1);\draw[te] (1,-1)--(2,-2);\draw[te] (3,0)--(5,-1);\draw[te] (5,-1)--(4,-2);\node[bn,minimum size=2.70mm,font=\tiny] at (3,0) {3};\node[bn,minimum size=2.70mm,font=\tiny] at (1,-1) {1};\node[bn,minimum size=2.70mm,font=\tiny] at (2,-2) {2};\node[bn,minimum size=2.70mm,font=\tiny] at (5,-1) {5};\node[bn,minimum size=2.70mm,font=\tiny] at (4,-2) {4};\end{tikzpicture}};
\node[vbox,fill=assocA] (v21) at (-3.3120,11.7114) {\begin{tikzpicture}[x=0.340cm,y=0.300cm]\path (0.50,0.55) rectangle (5.50,-2.55);\draw[te] (3,0)--(2,-1);\draw[te] (2,-1)--(1,-2);\draw[te] (3,0)--(4,-1);\draw[te] (4,-1)--(5,-2);\node[bn,minimum size=2.70mm,font=\tiny] at (3,0) {3};\node[bn,minimum size=2.70mm,font=\tiny] at (2,-1) {2};\node[bn,minimum size=2.70mm,font=\tiny] at (1,-2) {1};\node[bn,minimum size=2.70mm,font=\tiny] at (4,-1) {4};\node[bn,minimum size=2.70mm,font=\tiny] at (5,-2) {5};\end{tikzpicture}};
\node[vbox,fill=assocC] (v22) at (-1.6451,9.2456) {\begin{tikzpicture}[x=0.340cm,y=0.300cm]\path (0.50,0.55) rectangle (5.50,-2.55);\draw[te] (3,0)--(2,-1);\draw[te] (2,-1)--(1,-2);\draw[te] (3,0)--(5,-1);\draw[te] (5,-1)--(4,-2);\node[bn,minimum size=2.70mm,font=\tiny] at (3,0) {3};\node[bn,minimum size=2.70mm,font=\tiny] at (2,-1) {2};\node[bn,minimum size=2.70mm,font=\tiny] at (1,-2) {1};\node[bn,minimum size=2.70mm,font=\tiny] at (5,-1) {5};\node[bn,minimum size=2.70mm,font=\tiny] at (4,-2) {4};\end{tikzpicture}};
\node[vbox,fill=assocA] (v23) at (-6.8246,-1.7975) {\begin{tikzpicture}[x=0.340cm,y=0.300cm]\path (0.50,0.55) rectangle (5.50,-3.55);\draw[te] (4,0)--(1,-1);\draw[te] (1,-1)--(2,-2);\draw[te] (2,-2)--(3,-3);\draw[te] (4,0)--(5,-1);\node[bn,minimum size=2.70mm,font=\tiny] at (4,0) {4};\node[bn,minimum size=2.70mm,font=\tiny] at (1,-1) {1};\node[bn,minimum size=2.70mm,font=\tiny] at (2,-2) {2};\node[bn,minimum size=2.70mm,font=\tiny] at (3,-3) {3};\node[bn,minimum size=2.70mm,font=\tiny] at (5,-1) {5};\end{tikzpicture}};
\node[vbox,fill=assocC] (v24) at (-9.3468,-3.0777) {\begin{tikzpicture}[x=0.340cm,y=0.300cm]\path (0.50,0.55) rectangle (5.50,-3.55);\draw[te] (4,0)--(1,-1);\draw[te] (1,-1)--(3,-2);\draw[te] (3,-2)--(2,-3);\draw[te] (4,0)--(5,-1);\node[bn,minimum size=2.70mm,font=\tiny] at (4,0) {4};\node[bn,minimum size=2.70mm,font=\tiny] at (1,-1) {1};\node[bn,minimum size=2.70mm,font=\tiny] at (3,-2) {3};\node[bn,minimum size=2.70mm,font=\tiny] at (2,-3) {2};\node[bn,minimum size=2.70mm,font=\tiny] at (5,-1) {5};\end{tikzpicture}};
\node[vbox,fill=assocD] (v25) at (-6.3896,1.1462) {\begin{tikzpicture}[x=0.340cm,y=0.300cm]\path (0.50,0.55) rectangle (5.50,-2.55);\draw[te] (4,0)--(2,-1);\draw[te] (2,-1)--(1,-2);\draw[te] (2,-1)--(3,-2);\draw[te] (4,0)--(5,-1);\node[bn,minimum size=2.70mm,font=\tiny] at (4,0) {4};\node[bn,minimum size=2.70mm,font=\tiny] at (2,-1) {2};\node[bn,minimum size=2.70mm,font=\tiny] at (1,-2) {1};\node[bn,minimum size=2.70mm,font=\tiny] at (3,-2) {3};\node[bn,minimum size=2.70mm,font=\tiny] at (5,-1) {5};\end{tikzpicture}};
\node[vbox,fill=assocA] (v26) at (-11.4304,5.1725) {\begin{tikzpicture}[x=0.340cm,y=0.300cm]\path (0.50,0.55) rectangle (5.50,-3.55);\draw[te] (4,0)--(3,-1);\draw[te] (3,-1)--(1,-2);\draw[te] (1,-2)--(2,-3);\draw[te] (4,0)--(5,-1);\node[bn,minimum size=2.70mm,font=\tiny] at (4,0) {4};\node[bn,minimum size=2.70mm,font=\tiny] at (3,-1) {3};\node[bn,minimum size=2.70mm,font=\tiny] at (1,-2) {1};\node[bn,minimum size=2.70mm,font=\tiny] at (2,-3) {2};\node[bn,minimum size=2.70mm,font=\tiny] at (5,-1) {5};\end{tikzpicture}};
\node[vbox,fill=assocC] (v27) at (-8.4632,4.9384) {\begin{tikzpicture}[x=0.340cm,y=0.300cm]\path (0.50,0.55) rectangle (5.50,-3.55);\draw[te] (4,0)--(3,-1);\draw[te] (3,-1)--(2,-2);\draw[te] (2,-2)--(1,-3);\draw[te] (4,0)--(5,-1);\node[bn,minimum size=2.70mm,font=\tiny] at (4,0) {4};\node[bn,minimum size=2.70mm,font=\tiny] at (3,-1) {3};\node[bn,minimum size=2.70mm,font=\tiny] at (2,-2) {2};\node[bn,minimum size=2.70mm,font=\tiny] at (1,-3) {1};\node[bn,minimum size=2.70mm,font=\tiny] at (5,-1) {5};\end{tikzpicture}};
\node[vbox,fill=assocA] (v28) at (2.5720,-5.6701) {\begin{tikzpicture}[x=0.340cm,y=0.300cm]\path (0.50,0.55) rectangle (5.50,-4.55);\draw[te] (5,0)--(1,-1);\draw[te] (1,-1)--(2,-2);\draw[te] (2,-2)--(3,-3);\draw[te] (3,-3)--(4,-4);\node[bn,minimum size=2.70mm,font=\tiny] at (5,0) {5};\node[bn,minimum size=2.70mm,font=\tiny] at (1,-1) {1};\node[bn,minimum size=2.70mm,font=\tiny] at (2,-2) {2};\node[bn,minimum size=2.70mm,font=\tiny] at (3,-3) {3};\node[bn,minimum size=2.70mm,font=\tiny] at (4,-4) {4};\end{tikzpicture}};
\node[vbox,fill=assocC] (v29) at (4.3813,-7.8442) {\begin{tikzpicture}[x=0.340cm,y=0.300cm]\path (0.50,0.55) rectangle (5.50,-4.55);\draw[te] (5,0)--(1,-1);\draw[te] (1,-1)--(2,-2);\draw[te] (2,-2)--(4,-3);\draw[te] (4,-3)--(3,-4);\node[bn,minimum size=2.70mm,font=\tiny] at (5,0) {5};\node[bn,minimum size=2.70mm,font=\tiny] at (1,-1) {1};\node[bn,minimum size=2.70mm,font=\tiny] at (2,-2) {2};\node[bn,minimum size=2.70mm,font=\tiny] at (4,-3) {4};\node[bn,minimum size=2.70mm,font=\tiny] at (3,-4) {3};\end{tikzpicture}};
\node[vbox,fill=assocD] (v30) at (-0.3947,-5.9010) {\begin{tikzpicture}[x=0.340cm,y=0.300cm]\path (0.50,0.55) rectangle (5.50,-3.55);\draw[te] (5,0)--(1,-1);\draw[te] (1,-1)--(3,-2);\draw[te] (3,-2)--(2,-3);\draw[te] (3,-2)--(4,-3);\node[bn,minimum size=2.70mm,font=\tiny] at (5,0) {5};\node[bn,minimum size=2.70mm,font=\tiny] at (1,-1) {1};\node[bn,minimum size=2.70mm,font=\tiny] at (3,-2) {3};\node[bn,minimum size=2.70mm,font=\tiny] at (2,-3) {2};\node[bn,minimum size=2.70mm,font=\tiny] at (4,-3) {4};\end{tikzpicture}};
\node[vbox,fill=assocA] (v31) at (-3.1984,-11.7114) {\begin{tikzpicture}[x=0.340cm,y=0.300cm]\path (0.50,0.55) rectangle (5.50,-4.55);\draw[te] (5,0)--(1,-1);\draw[te] (1,-1)--(4,-2);\draw[te] (4,-2)--(2,-3);\draw[te] (2,-3)--(3,-4);\node[bn,minimum size=2.70mm,font=\tiny] at (5,0) {5};\node[bn,minimum size=2.70mm,font=\tiny] at (1,-1) {1};\node[bn,minimum size=2.70mm,font=\tiny] at (4,-2) {4};\node[bn,minimum size=2.70mm,font=\tiny] at (2,-3) {2};\node[bn,minimum size=2.70mm,font=\tiny] at (3,-4) {3};\end{tikzpicture}};
\node[vbox,fill=assocC] (v32) at (-3.6304,-8.7665) {\begin{tikzpicture}[x=0.340cm,y=0.300cm]\path (0.50,0.55) rectangle (5.50,-4.55);\draw[te] (5,0)--(1,-1);\draw[te] (1,-1)--(4,-2);\draw[te] (4,-2)--(3,-3);\draw[te] (3,-3)--(2,-4);\node[bn,minimum size=2.70mm,font=\tiny] at (5,0) {5};\node[bn,minimum size=2.70mm,font=\tiny] at (1,-1) {1};\node[bn,minimum size=2.70mm,font=\tiny] at (4,-2) {4};\node[bn,minimum size=2.70mm,font=\tiny] at (3,-3) {3};\node[bn,minimum size=2.70mm,font=\tiny] at (2,-4) {2};\end{tikzpicture}};
\node[vbox,fill=assocC] (v33) at (2.7487,-1.2925) {\begin{tikzpicture}[x=0.340cm,y=0.300cm]\path (0.50,0.55) rectangle (5.50,-3.55);\draw[te] (5,0)--(2,-1);\draw[te] (2,-1)--(1,-2);\draw[te] (2,-1)--(3,-2);\draw[te] (3,-2)--(4,-3);\node[bn,minimum size=2.70mm,font=\tiny] at (5,0) {5};\node[bn,minimum size=2.70mm,font=\tiny] at (2,-1) {2};\node[bn,minimum size=2.70mm,font=\tiny] at (1,-2) {1};\node[bn,minimum size=2.70mm,font=\tiny] at (3,-2) {3};\node[bn,minimum size=2.70mm,font=\tiny] at (4,-3) {4};\end{tikzpicture}};
\node[vbox,fill=assocD] (v34) at (4.1585,-3.5278) {\begin{tikzpicture}[x=0.340cm,y=0.300cm]\path (0.50,0.55) rectangle (5.50,-3.55);\draw[te] (5,0)--(2,-1);\draw[te] (2,-1)--(1,-2);\draw[te] (2,-1)--(4,-2);\draw[te] (4,-2)--(3,-3);\node[bn,minimum size=2.70mm,font=\tiny] at (5,0) {5};\node[bn,minimum size=2.70mm,font=\tiny] at (2,-1) {2};\node[bn,minimum size=2.70mm,font=\tiny] at (1,-2) {1};\node[bn,minimum size=2.70mm,font=\tiny] at (4,-2) {4};\node[bn,minimum size=2.70mm,font=\tiny] at (3,-3) {3};\end{tikzpicture}};
\node[vbox,fill=assocC] (v35) at (-3.0168,2.6344) {\begin{tikzpicture}[x=0.340cm,y=0.300cm]\path (0.50,0.55) rectangle (5.50,-3.55);\draw[te] (5,0)--(3,-1);\draw[te] (3,-1)--(1,-2);\draw[te] (1,-2)--(2,-3);\draw[te] (3,-1)--(4,-2);\node[bn,minimum size=2.70mm,font=\tiny] at (5,0) {5};\node[bn,minimum size=2.70mm,font=\tiny] at (3,-1) {3};\node[bn,minimum size=2.70mm,font=\tiny] at (1,-2) {1};\node[bn,minimum size=2.70mm,font=\tiny] at (2,-3) {2};\node[bn,minimum size=2.70mm,font=\tiny] at (4,-2) {4};\end{tikzpicture}};
\node[vbox,fill=assocD] (v36) at (-3.3171,5.2600) {\begin{tikzpicture}[x=0.340cm,y=0.300cm]\path (0.50,0.55) rectangle (5.50,-3.55);\draw[te] (5,0)--(3,-1);\draw[te] (3,-1)--(2,-2);\draw[te] (2,-2)--(1,-3);\draw[te] (3,-1)--(4,-2);\node[bn,minimum size=2.70mm,font=\tiny] at (5,0) {5};\node[bn,minimum size=2.70mm,font=\tiny] at (3,-1) {3};\node[bn,minimum size=2.70mm,font=\tiny] at (2,-2) {2};\node[bn,minimum size=2.70mm,font=\tiny] at (1,-3) {1};\node[bn,minimum size=2.70mm,font=\tiny] at (4,-2) {4};\end{tikzpicture}};
\node[vbox,fill=assocD] (v37) at (-7.6657,-5.3434) {\begin{tikzpicture}[x=0.340cm,y=0.300cm]\path (0.50,0.55) rectangle (5.50,-4.55);\draw[te] (5,0)--(4,-1);\draw[te] (4,-1)--(1,-2);\draw[te] (1,-2)--(2,-3);\draw[te] (2,-3)--(3,-4);\node[bn,minimum size=2.70mm,font=\tiny] at (5,0) {5};\node[bn,minimum size=2.70mm,font=\tiny] at (4,-1) {4};\node[bn,minimum size=2.70mm,font=\tiny] at (1,-2) {1};\node[bn,minimum size=2.70mm,font=\tiny] at (2,-3) {2};\node[bn,minimum size=2.70mm,font=\tiny] at (3,-4) {3};\end{tikzpicture}};
\node[vbox,fill=assocA] (v38) at (-11.3799,-5.2516) {\begin{tikzpicture}[x=0.340cm,y=0.300cm]\path (0.50,0.55) rectangle (5.50,-4.55);\draw[te] (5,0)--(4,-1);\draw[te] (4,-1)--(1,-2);\draw[te] (1,-2)--(3,-3);\draw[te] (3,-3)--(2,-4);\node[bn,minimum size=2.70mm,font=\tiny] at (5,0) {5};\node[bn,minimum size=2.70mm,font=\tiny] at (4,-1) {4};\node[bn,minimum size=2.70mm,font=\tiny] at (1,-2) {1};\node[bn,minimum size=2.70mm,font=\tiny] at (3,-3) {3};\node[bn,minimum size=2.70mm,font=\tiny] at (2,-4) {2};\end{tikzpicture}};
\node[vbox,fill=assocC] (v39) at (-4.1496,-0.2560) {\begin{tikzpicture}[x=0.340cm,y=0.300cm]\path (0.50,0.55) rectangle (5.50,-3.55);\draw[te] (5,0)--(4,-1);\draw[te] (4,-1)--(2,-2);\draw[te] (2,-2)--(1,-3);\draw[te] (2,-2)--(3,-3);\node[bn,minimum size=2.70mm,font=\tiny] at (5,0) {5};\node[bn,minimum size=2.70mm,font=\tiny] at (4,-1) {4};\node[bn,minimum size=2.70mm,font=\tiny] at (2,-2) {2};\node[bn,minimum size=2.70mm,font=\tiny] at (1,-3) {1};\node[bn,minimum size=2.70mm,font=\tiny] at (3,-3) {3};\end{tikzpicture}};
\node[vbox,fill=assocD] (v40) at (-9.1865,2.2114) {\begin{tikzpicture}[x=0.340cm,y=0.300cm]\path (0.50,0.55) rectangle (5.50,-4.55);\draw[te] (5,0)--(4,-1);\draw[te] (4,-1)--(3,-2);\draw[te] (3,-2)--(1,-3);\draw[te] (1,-3)--(2,-4);\node[bn,minimum size=2.70mm,font=\tiny] at (5,0) {5};\node[bn,minimum size=2.70mm,font=\tiny] at (4,-1) {4};\node[bn,minimum size=2.70mm,font=\tiny] at (3,-2) {3};\node[bn,minimum size=2.70mm,font=\tiny] at (1,-3) {1};\node[bn,minimum size=2.70mm,font=\tiny] at (2,-4) {2};\end{tikzpicture}};
\node[vbox,fill=assocA] (v41) at (-5.8898,3.7646) {\begin{tikzpicture}[x=0.340cm,y=0.300cm]\path (0.50,0.55) rectangle (5.50,-4.55);\draw[te] (5,0)--(4,-1);\draw[te] (4,-1)--(3,-2);\draw[te] (3,-2)--(2,-3);\draw[te] (2,-3)--(1,-4);\node[bn,minimum size=2.70mm,font=\tiny] at (5,0) {5};\node[bn,minimum size=2.70mm,font=\tiny] at (4,-1) {4};\node[bn,minimum size=2.70mm,font=\tiny] at (3,-2) {3};\node[bn,minimum size=2.70mm,font=\tiny] at (2,-3) {2};\node[bn,minimum size=2.70mm,font=\tiny] at (1,-4) {1};\end{tikzpicture}};
\begin{pgfonlayer}{bg}
\draw[gedge] (v0)--(v1);
\draw[gedge] (v0)--(v2);
\draw[gedge] (v0)--(v5);
\draw[gedge] (v0)--(v14);
\draw[gedge] (v1)--(v3);
\draw[gedge] (v1)--(v6);
\draw[gedge] (v1)--(v15);
\draw[gedge] (v2)--(v4);
\draw[gedge] (v2)--(v7);
\draw[gedge] (v2)--(v16);
\draw[gedge] (v3)--(v4);
\draw[gedge] (v3)--(v9);
\draw[gedge] (v3)--(v17);
\draw[gedge] (v4)--(v10);
\draw[gedge] (v4)--(v18);
\draw[gedge] (v5)--(v6);
\draw[gedge] (v5)--(v8);
\draw[gedge] (v5)--(v19);
\draw[gedge] (v6)--(v11);
\draw[gedge] (v6)--(v20);
\draw[gedge] (v7)--(v8);
\draw[gedge] (v7)--(v12);
\draw[gedge] (v7)--(v23);
\draw[gedge] (v8)--(v13);
\draw[gedge] (v8)--(v24);
\draw[gedge] (v9)--(v10);
\draw[gedge] (v9)--(v11);
\draw[gedge] (v9)--(v28);
\draw[gedge] (v10)--(v12);
\draw[gedge] (v10)--(v29);
\draw[gedge] (v11)--(v13);
\draw[gedge] (v11)--(v30);
\draw[gedge] (v12)--(v13);
\draw[gedge] (v12)--(v31);
\draw[gedge] (v13)--(v32);
\draw[gedge] (v14)--(v15);
\draw[gedge] (v14)--(v16);
\draw[gedge] (v14)--(v21);
\draw[gedge] (v15)--(v17);
\draw[gedge] (v15)--(v22);
\draw[gedge] (v16)--(v18);
\draw[gedge] (v16)--(v25);
\draw[gedge] (v17)--(v18);
\draw[gedge] (v17)--(v33);
\draw[gedge] (v18)--(v34);
\draw[gedge] (v19)--(v20);
\draw[gedge] (v19)--(v21);
\draw[gedge] (v19)--(v26);
\draw[gedge] (v20)--(v22);
\draw[gedge] (v20)--(v35);
\draw[gedge] (v21)--(v22);
\draw[gedge] (v21)--(v27);
\draw[gedge] (v22)--(v36);
\draw[gedge] (v23)--(v24);
\draw[gedge] (v23)--(v25);
\draw[gedge] (v23)--(v37);
\draw[gedge] (v24)--(v26);
\draw[gedge] (v24)--(v38);
\draw[gedge] (v25)--(v27);
\draw[gedge] (v25)--(v39);
\draw[gedge] (v26)--(v27);
\draw[gedge] (v26)--(v40);
\draw[gedge] (v27)--(v41);
\draw[gedge] (v28)--(v29);
\draw[gedge] (v28)--(v30);
\draw[gedge] (v28)--(v33);
\draw[gedge] (v29)--(v31);
\draw[gedge] (v29)--(v34);
\draw[gedge] (v30)--(v32);
\draw[gedge] (v30)--(v35);
\draw[gedge] (v31)--(v32);
\draw[gedge] (v31)--(v37);
\draw[gedge] (v32)--(v38);
\draw[gedge] (v33)--(v34);
\draw[gedge] (v33)--(v36);
\draw[gedge] (v34)--(v39);
\draw[gedge] (v35)--(v36);
\draw[gedge] (v35)--(v40);
\draw[gedge] (v36)--(v41);
\draw[gedge] (v37)--(v38);
\draw[gedge] (v37)--(v39);
\draw[gedge] (v38)--(v40);
\draw[gedge] (v39)--(v41);
\draw[gedge] (v40)--(v41);
\end{pgfonlayer}
\node[vbox,fill=assocA] at (4.2565,4.3527) {\begin{tikzpicture}[x=0.340cm,y=0.300cm]\path (0.50,0.55) rectangle (5.50,-4.55);\draw[te] (1,0)--(2,-1);\draw[te] (2,-1)--(3,-2);\draw[te] (3,-2)--(4,-3);\draw[te] (4,-3)--(5,-4);\node[bn,minimum size=2.70mm,font=\tiny] at (1,0) {1};\node[bn,minimum size=2.70mm,font=\tiny] at (2,-1) {2};\node[bn,minimum size=2.70mm,font=\tiny] at (3,-2) {3};\node[bn,minimum size=2.70mm,font=\tiny] at (4,-3) {4};\node[bn,minimum size=2.70mm,font=\tiny] at (5,-4) {5};\end{tikzpicture}};
\node[vbox,fill=assocC] at (5.9735,6.6004) {\begin{tikzpicture}[x=0.340cm,y=0.300cm]\path (0.50,0.55) rectangle (5.50,-4.55);\draw[te] (1,0)--(2,-1);\draw[te] (2,-1)--(3,-2);\draw[te] (3,-2)--(5,-3);\draw[te] (5,-3)--(4,-4);\node[bn,minimum size=2.70mm,font=\tiny] at (1,0) {1};\node[bn,minimum size=2.70mm,font=\tiny] at (2,-1) {2};\node[bn,minimum size=2.70mm,font=\tiny] at (3,-2) {3};\node[bn,minimum size=2.70mm,font=\tiny] at (5,-3) {5};\node[bn,minimum size=2.70mm,font=\tiny] at (4,-4) {4};\end{tikzpicture}};
\node[vbox,fill=assocD] at (5.1418,1.5118) {\begin{tikzpicture}[x=0.340cm,y=0.300cm]\path (0.50,0.55) rectangle (5.50,-3.55);\draw[te] (1,0)--(2,-1);\draw[te] (2,-1)--(4,-2);\draw[te] (4,-2)--(3,-3);\draw[te] (4,-2)--(5,-3);\node[bn,minimum size=2.70mm,font=\tiny] at (1,0) {1};\node[bn,minimum size=2.70mm,font=\tiny] at (2,-1) {2};\node[bn,minimum size=2.70mm,font=\tiny] at (4,-2) {4};\node[bn,minimum size=2.70mm,font=\tiny] at (3,-3) {3};\node[bn,minimum size=2.70mm,font=\tiny] at (5,-3) {5};\end{tikzpicture}};
\node[vbox,fill=assocA] at (11.4304,0.0713) {\begin{tikzpicture}[x=0.340cm,y=0.300cm]\path (0.50,0.55) rectangle (5.50,-4.55);\draw[te] (1,0)--(2,-1);\draw[te] (2,-1)--(5,-2);\draw[te] (5,-2)--(3,-3);\draw[te] (3,-3)--(4,-4);\node[bn,minimum size=2.70mm,font=\tiny] at (1,0) {1};\node[bn,minimum size=2.70mm,font=\tiny] at (2,-1) {2};\node[bn,minimum size=2.70mm,font=\tiny] at (5,-2) {5};\node[bn,minimum size=2.70mm,font=\tiny] at (3,-3) {3};\node[bn,minimum size=2.70mm,font=\tiny] at (4,-4) {4};\end{tikzpicture}};
\node[vbox,fill=assocC] at (8.6555,-1.0052) {\begin{tikzpicture}[x=0.340cm,y=0.300cm]\path (0.50,0.55) rectangle (5.50,-4.55);\draw[te] (1,0)--(2,-1);\draw[te] (2,-1)--(5,-2);\draw[te] (5,-2)--(4,-3);\draw[te] (4,-3)--(3,-4);\node[bn,minimum size=2.70mm,font=\tiny] at (1,0) {1};\node[bn,minimum size=2.70mm,font=\tiny] at (2,-1) {2};\node[bn,minimum size=2.70mm,font=\tiny] at (5,-2) {5};\node[bn,minimum size=2.70mm,font=\tiny] at (4,-3) {4};\node[bn,minimum size=2.70mm,font=\tiny] at (3,-4) {3};\end{tikzpicture}};
\node[vbox,fill=assocC] at (-0.0506,3.5508) {\begin{tikzpicture}[x=0.340cm,y=0.300cm]\path (0.50,0.55) rectangle (5.50,-3.55);\draw[te] (1,0)--(3,-1);\draw[te] (3,-1)--(2,-2);\draw[te] (3,-1)--(4,-2);\draw[te] (4,-2)--(5,-3);\node[bn,minimum size=2.70mm,font=\tiny] at (1,0) {1};\node[bn,minimum size=2.70mm,font=\tiny] at (3,-1) {3};\node[bn,minimum size=2.70mm,font=\tiny] at (2,-2) {2};\node[bn,minimum size=2.70mm,font=\tiny] at (4,-2) {4};\node[bn,minimum size=2.70mm,font=\tiny] at (5,-3) {5};\end{tikzpicture}};
\node[vbox,fill=assocD] at (1.8149,5.4227) {\begin{tikzpicture}[x=0.340cm,y=0.300cm]\path (0.50,0.55) rectangle (5.50,-3.55);\draw[te] (1,0)--(3,-1);\draw[te] (3,-1)--(2,-2);\draw[te] (3,-1)--(5,-2);\draw[te] (5,-2)--(4,-3);\node[bn,minimum size=2.70mm,font=\tiny] at (1,0) {1};\node[bn,minimum size=2.70mm,font=\tiny] at (3,-1) {3};\node[bn,minimum size=2.70mm,font=\tiny] at (2,-2) {2};\node[bn,minimum size=2.70mm,font=\tiny] at (5,-2) {5};\node[bn,minimum size=2.70mm,font=\tiny] at (4,-3) {4};\end{tikzpicture}};
\node[vbox,fill=assocC] at (-2.5961,-2.9439) {\begin{tikzpicture}[x=0.340cm,y=0.300cm]\path (0.50,0.55) rectangle (5.50,-3.55);\draw[te] (1,0)--(4,-1);\draw[te] (4,-1)--(2,-2);\draw[te] (2,-2)--(3,-3);\draw[te] (4,-1)--(5,-2);\node[bn,minimum size=2.70mm,font=\tiny] at (1,0) {1};\node[bn,minimum size=2.70mm,font=\tiny] at (4,-1) {4};\node[bn,minimum size=2.70mm,font=\tiny] at (2,-2) {2};\node[bn,minimum size=2.70mm,font=\tiny] at (3,-3) {3};\node[bn,minimum size=2.70mm,font=\tiny] at (5,-2) {5};\end{tikzpicture}};
\node[vbox,fill=assocD] at (-5.0890,-3.8209) {\begin{tikzpicture}[x=0.340cm,y=0.300cm]\path (0.50,0.55) rectangle (5.50,-3.55);\draw[te] (1,0)--(4,-1);\draw[te] (4,-1)--(3,-2);\draw[te] (3,-2)--(2,-3);\draw[te] (4,-1)--(5,-2);\node[bn,minimum size=2.70mm,font=\tiny] at (1,0) {1};\node[bn,minimum size=2.70mm,font=\tiny] at (4,-1) {4};\node[bn,minimum size=2.70mm,font=\tiny] at (3,-2) {3};\node[bn,minimum size=2.70mm,font=\tiny] at (2,-3) {2};\node[bn,minimum size=2.70mm,font=\tiny] at (5,-2) {5};\end{tikzpicture}};
\node[vbox,fill=assocD] at (6.2162,-5.7010) {\begin{tikzpicture}[x=0.340cm,y=0.300cm]\path (0.50,0.55) rectangle (5.50,-4.55);\draw[te] (1,0)--(5,-1);\draw[te] (5,-1)--(2,-2);\draw[te] (2,-2)--(3,-3);\draw[te] (3,-3)--(4,-4);\node[bn,minimum size=2.70mm,font=\tiny] at (1,0) {1};\node[bn,minimum size=2.70mm,font=\tiny] at (5,-1) {5};\node[bn,minimum size=2.70mm,font=\tiny] at (2,-2) {2};\node[bn,minimum size=2.70mm,font=\tiny] at (3,-3) {3};\node[bn,minimum size=2.70mm,font=\tiny] at (4,-4) {4};\end{tikzpicture}};
\node[vbox,fill=assocA] at (6.9531,-9.3425) {\begin{tikzpicture}[x=0.340cm,y=0.300cm]\path (0.50,0.55) rectangle (5.50,-4.55);\draw[te] (1,0)--(5,-1);\draw[te] (5,-1)--(2,-2);\draw[te] (2,-2)--(4,-3);\draw[te] (4,-3)--(3,-4);\node[bn,minimum size=2.70mm,font=\tiny] at (1,0) {1};\node[bn,minimum size=2.70mm,font=\tiny] at (5,-1) {5};\node[bn,minimum size=2.70mm,font=\tiny] at (2,-2) {2};\node[bn,minimum size=2.70mm,font=\tiny] at (4,-3) {4};\node[bn,minimum size=2.70mm,font=\tiny] at (3,-4) {3};\end{tikzpicture}};
\node[vbox,fill=assocC] at (0.4739,-3.4052) {\begin{tikzpicture}[x=0.340cm,y=0.300cm]\path (0.50,0.55) rectangle (5.50,-3.55);\draw[te] (1,0)--(5,-1);\draw[te] (5,-1)--(3,-2);\draw[te] (3,-2)--(2,-3);\draw[te] (3,-2)--(4,-3);\node[bn,minimum size=2.70mm,font=\tiny] at (1,0) {1};\node[bn,minimum size=2.70mm,font=\tiny] at (5,-1) {5};\node[bn,minimum size=2.70mm,font=\tiny] at (3,-2) {3};\node[bn,minimum size=2.70mm,font=\tiny] at (2,-3) {2};\node[bn,minimum size=2.70mm,font=\tiny] at (4,-3) {4};\end{tikzpicture}};
\node[vbox,fill=assocD] at (-0.8109,-8.8648) {\begin{tikzpicture}[x=0.340cm,y=0.300cm]\path (0.50,0.55) rectangle (5.50,-4.55);\draw[te] (1,0)--(5,-1);\draw[te] (5,-1)--(4,-2);\draw[te] (4,-2)--(2,-3);\draw[te] (2,-3)--(3,-4);\node[bn,minimum size=2.70mm,font=\tiny] at (1,0) {1};\node[bn,minimum size=2.70mm,font=\tiny] at (5,-1) {5};\node[bn,minimum size=2.70mm,font=\tiny] at (4,-2) {4};\node[bn,minimum size=2.70mm,font=\tiny] at (2,-3) {2};\node[bn,minimum size=2.70mm,font=\tiny] at (3,-4) {3};\end{tikzpicture}};
\node[vbox,fill=assocA] at (-3.0587,-5.9964) {\begin{tikzpicture}[x=0.340cm,y=0.300cm]\path (0.50,0.55) rectangle (5.50,-4.55);\draw[te] (1,0)--(5,-1);\draw[te] (5,-1)--(4,-2);\draw[te] (4,-2)--(3,-3);\draw[te] (3,-3)--(2,-4);\node[bn,minimum size=2.70mm,font=\tiny] at (1,0) {1};\node[bn,minimum size=2.70mm,font=\tiny] at (5,-1) {5};\node[bn,minimum size=2.70mm,font=\tiny] at (4,-2) {4};\node[bn,minimum size=2.70mm,font=\tiny] at (3,-3) {3};\node[bn,minimum size=2.70mm,font=\tiny] at (2,-4) {2};\end{tikzpicture}};
\node[vbox,fill=assocD] at (3.4758,7.9123) {\begin{tikzpicture}[x=0.340cm,y=0.300cm]\path (0.50,0.55) rectangle (5.50,-3.55);\draw[te] (2,0)--(1,-1);\draw[te] (2,0)--(3,-1);\draw[te] (3,-1)--(4,-2);\draw[te] (4,-2)--(5,-3);\node[bn,minimum size=2.70mm,font=\tiny] at (2,0) {2};\node[bn,minimum size=2.70mm,font=\tiny] at (1,-1) {1};\node[bn,minimum size=2.70mm,font=\tiny] at (3,-1) {3};\node[bn,minimum size=2.70mm,font=\tiny] at (4,-2) {4};\node[bn,minimum size=2.70mm,font=\tiny] at (5,-3) {5};\end{tikzpicture}};
\node[vbox,fill=assocA] at (6.8620,9.4411) {\begin{tikzpicture}[x=0.340cm,y=0.300cm]\path (0.50,0.55) rectangle (5.50,-3.55);\draw[te] (2,0)--(1,-1);\draw[te] (2,0)--(3,-1);\draw[te] (3,-1)--(5,-2);\draw[te] (5,-2)--(4,-3);\node[bn,minimum size=2.70mm,font=\tiny] at (2,0) {2};\node[bn,minimum size=2.70mm,font=\tiny] at (1,-1) {1};\node[bn,minimum size=2.70mm,font=\tiny] at (3,-1) {3};\node[bn,minimum size=2.70mm,font=\tiny] at (5,-2) {5};\node[bn,minimum size=2.70mm,font=\tiny] at (4,-3) {4};\end{tikzpicture}};
\node[vbox,fill=assocC] at (2.5153,1.8032) {\begin{tikzpicture}[x=0.340cm,y=0.300cm]\path (0.50,0.55) rectangle (5.50,-2.55);\draw[te] (2,0)--(1,-1);\draw[te] (2,0)--(4,-1);\draw[te] (4,-1)--(3,-2);\draw[te] (4,-1)--(5,-2);\node[bn,minimum size=2.70mm,font=\tiny] at (2,0) {2};\node[bn,minimum size=2.70mm,font=\tiny] at (1,-1) {1};\node[bn,minimum size=2.70mm,font=\tiny] at (4,-1) {4};\node[bn,minimum size=2.70mm,font=\tiny] at (3,-2) {3};\node[bn,minimum size=2.70mm,font=\tiny] at (5,-2) {5};\end{tikzpicture}};
\node[vbox,fill=assocD] at (8.1239,1.7655) {\begin{tikzpicture}[x=0.340cm,y=0.300cm]\path (0.50,0.55) rectangle (5.50,-3.55);\draw[te] (2,0)--(1,-1);\draw[te] (2,0)--(5,-1);\draw[te] (5,-1)--(3,-2);\draw[te] (3,-2)--(4,-3);\node[bn,minimum size=2.70mm,font=\tiny] at (2,0) {2};\node[bn,minimum size=2.70mm,font=\tiny] at (1,-1) {1};\node[bn,minimum size=2.70mm,font=\tiny] at (5,-1) {5};\node[bn,minimum size=2.70mm,font=\tiny] at (3,-2) {3};\node[bn,minimum size=2.70mm,font=\tiny] at (4,-3) {4};\end{tikzpicture}};
\node[vbox,fill=assocA] at (5.8276,-1.0643) {\begin{tikzpicture}[x=0.340cm,y=0.300cm]\path (0.50,0.55) rectangle (5.50,-3.55);\draw[te] (2,0)--(1,-1);\draw[te] (2,0)--(5,-1);\draw[te] (5,-1)--(4,-2);\draw[te] (4,-2)--(3,-3);\node[bn,minimum size=2.70mm,font=\tiny] at (2,0) {2};\node[bn,minimum size=2.70mm,font=\tiny] at (1,-1) {1};\node[bn,minimum size=2.70mm,font=\tiny] at (5,-1) {5};\node[bn,minimum size=2.70mm,font=\tiny] at (4,-2) {4};\node[bn,minimum size=2.70mm,font=\tiny] at (3,-3) {3};\end{tikzpicture}};
\node[vbox,fill=assocD] at (-4.2280,8.1108) {\begin{tikzpicture}[x=0.340cm,y=0.300cm]\path (0.50,0.55) rectangle (5.50,-2.55);\draw[te] (3,0)--(1,-1);\draw[te] (1,-1)--(2,-2);\draw[te] (3,0)--(4,-1);\draw[te] (4,-1)--(5,-2);\node[bn,minimum size=2.70mm,font=\tiny] at (3,0) {3};\node[bn,minimum size=2.70mm,font=\tiny] at (1,-1) {1};\node[bn,minimum size=2.70mm,font=\tiny] at (2,-2) {2};\node[bn,minimum size=2.70mm,font=\tiny] at (4,-1) {4};\node[bn,minimum size=2.70mm,font=\tiny] at (5,-2) {5};\end{tikzpicture}};
\node[vbox,fill=assocA] at (-0.9582,6.5017) {\begin{tikzpicture}[x=0.340cm,y=0.300cm]\path (0.50,0.55) rectangle (5.50,-2.55);\draw[te] (3,0)--(1,-1);\draw[te] (1,-1)--(2,-2);\draw[te] (3,0)--(5,-1);\draw[te] (5,-1)--(4,-2);\node[bn,minimum size=2.70mm,font=\tiny] at (3,0) {3};\node[bn,minimum size=2.70mm,font=\tiny] at (1,-1) {1};\node[bn,minimum size=2.70mm,font=\tiny] at (2,-2) {2};\node[bn,minimum size=2.70mm,font=\tiny] at (5,-1) {5};\node[bn,minimum size=2.70mm,font=\tiny] at (4,-2) {4};\end{tikzpicture}};
\node[vbox,fill=assocA] at (-3.3120,11.7114) {\begin{tikzpicture}[x=0.340cm,y=0.300cm]\path (0.50,0.55) rectangle (5.50,-2.55);\draw[te] (3,0)--(2,-1);\draw[te] (2,-1)--(1,-2);\draw[te] (3,0)--(4,-1);\draw[te] (4,-1)--(5,-2);\node[bn,minimum size=2.70mm,font=\tiny] at (3,0) {3};\node[bn,minimum size=2.70mm,font=\tiny] at (2,-1) {2};\node[bn,minimum size=2.70mm,font=\tiny] at (1,-2) {1};\node[bn,minimum size=2.70mm,font=\tiny] at (4,-1) {4};\node[bn,minimum size=2.70mm,font=\tiny] at (5,-2) {5};\end{tikzpicture}};
\node[vbox,fill=assocC] at (-1.6451,9.2456) {\begin{tikzpicture}[x=0.340cm,y=0.300cm]\path (0.50,0.55) rectangle (5.50,-2.55);\draw[te] (3,0)--(2,-1);\draw[te] (2,-1)--(1,-2);\draw[te] (3,0)--(5,-1);\draw[te] (5,-1)--(4,-2);\node[bn,minimum size=2.70mm,font=\tiny] at (3,0) {3};\node[bn,minimum size=2.70mm,font=\tiny] at (2,-1) {2};\node[bn,minimum size=2.70mm,font=\tiny] at (1,-2) {1};\node[bn,minimum size=2.70mm,font=\tiny] at (5,-1) {5};\node[bn,minimum size=2.70mm,font=\tiny] at (4,-2) {4};\end{tikzpicture}};
\node[vbox,fill=assocA] at (-6.8246,-1.7975) {\begin{tikzpicture}[x=0.340cm,y=0.300cm]\path (0.50,0.55) rectangle (5.50,-3.55);\draw[te] (4,0)--(1,-1);\draw[te] (1,-1)--(2,-2);\draw[te] (2,-2)--(3,-3);\draw[te] (4,0)--(5,-1);\node[bn,minimum size=2.70mm,font=\tiny] at (4,0) {4};\node[bn,minimum size=2.70mm,font=\tiny] at (1,-1) {1};\node[bn,minimum size=2.70mm,font=\tiny] at (2,-2) {2};\node[bn,minimum size=2.70mm,font=\tiny] at (3,-3) {3};\node[bn,minimum size=2.70mm,font=\tiny] at (5,-1) {5};\end{tikzpicture}};
\node[vbox,fill=assocC] at (-9.3468,-3.0777) {\begin{tikzpicture}[x=0.340cm,y=0.300cm]\path (0.50,0.55) rectangle (5.50,-3.55);\draw[te] (4,0)--(1,-1);\draw[te] (1,-1)--(3,-2);\draw[te] (3,-2)--(2,-3);\draw[te] (4,0)--(5,-1);\node[bn,minimum size=2.70mm,font=\tiny] at (4,0) {4};\node[bn,minimum size=2.70mm,font=\tiny] at (1,-1) {1};\node[bn,minimum size=2.70mm,font=\tiny] at (3,-2) {3};\node[bn,minimum size=2.70mm,font=\tiny] at (2,-3) {2};\node[bn,minimum size=2.70mm,font=\tiny] at (5,-1) {5};\end{tikzpicture}};
\node[vbox,fill=assocD] at (-6.3896,1.1462) {\begin{tikzpicture}[x=0.340cm,y=0.300cm]\path (0.50,0.55) rectangle (5.50,-2.55);\draw[te] (4,0)--(2,-1);\draw[te] (2,-1)--(1,-2);\draw[te] (2,-1)--(3,-2);\draw[te] (4,0)--(5,-1);\node[bn,minimum size=2.70mm,font=\tiny] at (4,0) {4};\node[bn,minimum size=2.70mm,font=\tiny] at (2,-1) {2};\node[bn,minimum size=2.70mm,font=\tiny] at (1,-2) {1};\node[bn,minimum size=2.70mm,font=\tiny] at (3,-2) {3};\node[bn,minimum size=2.70mm,font=\tiny] at (5,-1) {5};\end{tikzpicture}};
\node[vbox,fill=assocA] at (-11.4304,5.1725) {\begin{tikzpicture}[x=0.340cm,y=0.300cm]\path (0.50,0.55) rectangle (5.50,-3.55);\draw[te] (4,0)--(3,-1);\draw[te] (3,-1)--(1,-2);\draw[te] (1,-2)--(2,-3);\draw[te] (4,0)--(5,-1);\node[bn,minimum size=2.70mm,font=\tiny] at (4,0) {4};\node[bn,minimum size=2.70mm,font=\tiny] at (3,-1) {3};\node[bn,minimum size=2.70mm,font=\tiny] at (1,-2) {1};\node[bn,minimum size=2.70mm,font=\tiny] at (2,-3) {2};\node[bn,minimum size=2.70mm,font=\tiny] at (5,-1) {5};\end{tikzpicture}};
\node[vbox,fill=assocC] at (-8.4632,4.9384) {\begin{tikzpicture}[x=0.340cm,y=0.300cm]\path (0.50,0.55) rectangle (5.50,-3.55);\draw[te] (4,0)--(3,-1);\draw[te] (3,-1)--(2,-2);\draw[te] (2,-2)--(1,-3);\draw[te] (4,0)--(5,-1);\node[bn,minimum size=2.70mm,font=\tiny] at (4,0) {4};\node[bn,minimum size=2.70mm,font=\tiny] at (3,-1) {3};\node[bn,minimum size=2.70mm,font=\tiny] at (2,-2) {2};\node[bn,minimum size=2.70mm,font=\tiny] at (1,-3) {1};\node[bn,minimum size=2.70mm,font=\tiny] at (5,-1) {5};\end{tikzpicture}};
\node[vbox,fill=assocA] at (2.5720,-5.6701) {\begin{tikzpicture}[x=0.340cm,y=0.300cm]\path (0.50,0.55) rectangle (5.50,-4.55);\draw[te] (5,0)--(1,-1);\draw[te] (1,-1)--(2,-2);\draw[te] (2,-2)--(3,-3);\draw[te] (3,-3)--(4,-4);\node[bn,minimum size=2.70mm,font=\tiny] at (5,0) {5};\node[bn,minimum size=2.70mm,font=\tiny] at (1,-1) {1};\node[bn,minimum size=2.70mm,font=\tiny] at (2,-2) {2};\node[bn,minimum size=2.70mm,font=\tiny] at (3,-3) {3};\node[bn,minimum size=2.70mm,font=\tiny] at (4,-4) {4};\end{tikzpicture}};
\node[vbox,fill=assocC] at (4.3813,-7.8442) {\begin{tikzpicture}[x=0.340cm,y=0.300cm]\path (0.50,0.55) rectangle (5.50,-4.55);\draw[te] (5,0)--(1,-1);\draw[te] (1,-1)--(2,-2);\draw[te] (2,-2)--(4,-3);\draw[te] (4,-3)--(3,-4);\node[bn,minimum size=2.70mm,font=\tiny] at (5,0) {5};\node[bn,minimum size=2.70mm,font=\tiny] at (1,-1) {1};\node[bn,minimum size=2.70mm,font=\tiny] at (2,-2) {2};\node[bn,minimum size=2.70mm,font=\tiny] at (4,-3) {4};\node[bn,minimum size=2.70mm,font=\tiny] at (3,-4) {3};\end{tikzpicture}};
\node[vbox,fill=assocD] at (-0.3947,-5.9010) {\begin{tikzpicture}[x=0.340cm,y=0.300cm]\path (0.50,0.55) rectangle (5.50,-3.55);\draw[te] (5,0)--(1,-1);\draw[te] (1,-1)--(3,-2);\draw[te] (3,-2)--(2,-3);\draw[te] (3,-2)--(4,-3);\node[bn,minimum size=2.70mm,font=\tiny] at (5,0) {5};\node[bn,minimum size=2.70mm,font=\tiny] at (1,-1) {1};\node[bn,minimum size=2.70mm,font=\tiny] at (3,-2) {3};\node[bn,minimum size=2.70mm,font=\tiny] at (2,-3) {2};\node[bn,minimum size=2.70mm,font=\tiny] at (4,-3) {4};\end{tikzpicture}};
\node[vbox,fill=assocA] at (-3.1984,-11.7114) {\begin{tikzpicture}[x=0.340cm,y=0.300cm]\path (0.50,0.55) rectangle (5.50,-4.55);\draw[te] (5,0)--(1,-1);\draw[te] (1,-1)--(4,-2);\draw[te] (4,-2)--(2,-3);\draw[te] (2,-3)--(3,-4);\node[bn,minimum size=2.70mm,font=\tiny] at (5,0) {5};\node[bn,minimum size=2.70mm,font=\tiny] at (1,-1) {1};\node[bn,minimum size=2.70mm,font=\tiny] at (4,-2) {4};\node[bn,minimum size=2.70mm,font=\tiny] at (2,-3) {2};\node[bn,minimum size=2.70mm,font=\tiny] at (3,-4) {3};\end{tikzpicture}};
\node[vbox,fill=assocC] at (-3.6304,-8.7665) {\begin{tikzpicture}[x=0.340cm,y=0.300cm]\path (0.50,0.55) rectangle (5.50,-4.55);\draw[te] (5,0)--(1,-1);\draw[te] (1,-1)--(4,-2);\draw[te] (4,-2)--(3,-3);\draw[te] (3,-3)--(2,-4);\node[bn,minimum size=2.70mm,font=\tiny] at (5,0) {5};\node[bn,minimum size=2.70mm,font=\tiny] at (1,-1) {1};\node[bn,minimum size=2.70mm,font=\tiny] at (4,-2) {4};\node[bn,minimum size=2.70mm,font=\tiny] at (3,-3) {3};\node[bn,minimum size=2.70mm,font=\tiny] at (2,-4) {2};\end{tikzpicture}};
\node[vbox,fill=assocC] at (2.7487,-1.2925) {\begin{tikzpicture}[x=0.340cm,y=0.300cm]\path (0.50,0.55) rectangle (5.50,-3.55);\draw[te] (5,0)--(2,-1);\draw[te] (2,-1)--(1,-2);\draw[te] (2,-1)--(3,-2);\draw[te] (3,-2)--(4,-3);\node[bn,minimum size=2.70mm,font=\tiny] at (5,0) {5};\node[bn,minimum size=2.70mm,font=\tiny] at (2,-1) {2};\node[bn,minimum size=2.70mm,font=\tiny] at (1,-2) {1};\node[bn,minimum size=2.70mm,font=\tiny] at (3,-2) {3};\node[bn,minimum size=2.70mm,font=\tiny] at (4,-3) {4};\end{tikzpicture}};
\node[vbox,fill=assocD] at (4.1585,-3.5278) {\begin{tikzpicture}[x=0.340cm,y=0.300cm]\path (0.50,0.55) rectangle (5.50,-3.55);\draw[te] (5,0)--(2,-1);\draw[te] (2,-1)--(1,-2);\draw[te] (2,-1)--(4,-2);\draw[te] (4,-2)--(3,-3);\node[bn,minimum size=2.70mm,font=\tiny] at (5,0) {5};\node[bn,minimum size=2.70mm,font=\tiny] at (2,-1) {2};\node[bn,minimum size=2.70mm,font=\tiny] at (1,-2) {1};\node[bn,minimum size=2.70mm,font=\tiny] at (4,-2) {4};\node[bn,minimum size=2.70mm,font=\tiny] at (3,-3) {3};\end{tikzpicture}};
\node[vbox,fill=assocC] at (-3.0168,2.6344) {\begin{tikzpicture}[x=0.340cm,y=0.300cm]\path (0.50,0.55) rectangle (5.50,-3.55);\draw[te] (5,0)--(3,-1);\draw[te] (3,-1)--(1,-2);\draw[te] (1,-2)--(2,-3);\draw[te] (3,-1)--(4,-2);\node[bn,minimum size=2.70mm,font=\tiny] at (5,0) {5};\node[bn,minimum size=2.70mm,font=\tiny] at (3,-1) {3};\node[bn,minimum size=2.70mm,font=\tiny] at (1,-2) {1};\node[bn,minimum size=2.70mm,font=\tiny] at (2,-3) {2};\node[bn,minimum size=2.70mm,font=\tiny] at (4,-2) {4};\end{tikzpicture}};
\node[vbox,fill=assocD] at (-3.3171,5.2600) {\begin{tikzpicture}[x=0.340cm,y=0.300cm]\path (0.50,0.55) rectangle (5.50,-3.55);\draw[te] (5,0)--(3,-1);\draw[te] (3,-1)--(2,-2);\draw[te] (2,-2)--(1,-3);\draw[te] (3,-1)--(4,-2);\node[bn,minimum size=2.70mm,font=\tiny] at (5,0) {5};\node[bn,minimum size=2.70mm,font=\tiny] at (3,-1) {3};\node[bn,minimum size=2.70mm,font=\tiny] at (2,-2) {2};\node[bn,minimum size=2.70mm,font=\tiny] at (1,-3) {1};\node[bn,minimum size=2.70mm,font=\tiny] at (4,-2) {4};\end{tikzpicture}};
\node[vbox,fill=assocD] at (-7.6657,-5.3434) {\begin{tikzpicture}[x=0.340cm,y=0.300cm]\path (0.50,0.55) rectangle (5.50,-4.55);\draw[te] (5,0)--(4,-1);\draw[te] (4,-1)--(1,-2);\draw[te] (1,-2)--(2,-3);\draw[te] (2,-3)--(3,-4);\node[bn,minimum size=2.70mm,font=\tiny] at (5,0) {5};\node[bn,minimum size=2.70mm,font=\tiny] at (4,-1) {4};\node[bn,minimum size=2.70mm,font=\tiny] at (1,-2) {1};\node[bn,minimum size=2.70mm,font=\tiny] at (2,-3) {2};\node[bn,minimum size=2.70mm,font=\tiny] at (3,-4) {3};\end{tikzpicture}};
\node[vbox,fill=assocA] at (-11.3799,-5.2516) {\begin{tikzpicture}[x=0.340cm,y=0.300cm]\path (0.50,0.55) rectangle (5.50,-4.55);\draw[te] (5,0)--(4,-1);\draw[te] (4,-1)--(1,-2);\draw[te] (1,-2)--(3,-3);\draw[te] (3,-3)--(2,-4);\node[bn,minimum size=2.70mm,font=\tiny] at (5,0) {5};\node[bn,minimum size=2.70mm,font=\tiny] at (4,-1) {4};\node[bn,minimum size=2.70mm,font=\tiny] at (1,-2) {1};\node[bn,minimum size=2.70mm,font=\tiny] at (3,-3) {3};\node[bn,minimum size=2.70mm,font=\tiny] at (2,-4) {2};\end{tikzpicture}};
\node[vbox,fill=assocC] at (-4.1496,-0.2560) {\begin{tikzpicture}[x=0.340cm,y=0.300cm]\path (0.50,0.55) rectangle (5.50,-3.55);\draw[te] (5,0)--(4,-1);\draw[te] (4,-1)--(2,-2);\draw[te] (2,-2)--(1,-3);\draw[te] (2,-2)--(3,-3);\node[bn,minimum size=2.70mm,font=\tiny] at (5,0) {5};\node[bn,minimum size=2.70mm,font=\tiny] at (4,-1) {4};\node[bn,minimum size=2.70mm,font=\tiny] at (2,-2) {2};\node[bn,minimum size=2.70mm,font=\tiny] at (1,-3) {1};\node[bn,minimum size=2.70mm,font=\tiny] at (3,-3) {3};\end{tikzpicture}};
\node[vbox,fill=assocD] at (-9.1865,2.2114) {\begin{tikzpicture}[x=0.340cm,y=0.300cm]\path (0.50,0.55) rectangle (5.50,-4.55);\draw[te] (5,0)--(4,-1);\draw[te] (4,-1)--(3,-2);\draw[te] (3,-2)--(1,-3);\draw[te] (1,-3)--(2,-4);\node[bn,minimum size=2.70mm,font=\tiny] at (5,0) {5};\node[bn,minimum size=2.70mm,font=\tiny] at (4,-1) {4};\node[bn,minimum size=2.70mm,font=\tiny] at (3,-2) {3};\node[bn,minimum size=2.70mm,font=\tiny] at (1,-3) {1};\node[bn,minimum size=2.70mm,font=\tiny] at (2,-4) {2};\end{tikzpicture}};
\node[vbox,fill=assocA] at (-5.8898,3.7646) {\begin{tikzpicture}[x=0.340cm,y=0.300cm]\path (0.50,0.55) rectangle (5.50,-4.55);\draw[te] (5,0)--(4,-1);\draw[te] (4,-1)--(3,-2);\draw[te] (3,-2)--(2,-3);\draw[te] (2,-3)--(1,-4);\node[bn,minimum size=2.70mm,font=\tiny] at (5,0) {5};\node[bn,minimum size=2.70mm,font=\tiny] at (4,-1) {4};\node[bn,minimum size=2.70mm,font=\tiny] at (3,-2) {3};\node[bn,minimum size=2.70mm,font=\tiny] at (2,-3) {2};\node[bn,minimum size=2.70mm,font=\tiny] at (1,-4) {1};\end{tikzpicture}};
\node[rlab] at ($(v0)!0.500!(v1)$) {4-5};
\node[rlab] at ($(v0)!0.500!(v2)$) {3-4};
\node[rlab] at ($(v0)!0.500!(v5)$) {2-3};
\node[rlab] at ($(v0)!0.500!(v14)$) {1-2};
\node[rlab] at ($(v1)!0.500!(v3)$) {3-5};
\node[rlab] at ($(v1)!0.380!(v6)$) {2-3};
\node[rlab] at ($(v1)!0.500!(v15)$) {1-2};
\node[rlab] at ($(v2)!0.380!(v4)$) {4-5};
\node[rlab] at ($(v2)!0.500!(v7)$) {2-4};
\node[rlab] at ($(v2)!0.500!(v16)$) {1-2};
\node[rlab] at ($(v3)!0.500!(v4)$) {3-4};
\node[rlab] at ($(v3)!0.500!(v9)$) {2-5};
\node[rlab] at ($(v3)!0.500!(v17)$) {1-2};
\node[rlab] at ($(v4)!0.500!(v10)$) {2-5};
\node[rlab] at ($(v4)!0.500!(v18)$) {1-2};
\node[rlab] at ($(v5)!0.500!(v6)$) {4-5};
\node[rlab] at ($(v5)!0.500!(v8)$) {3-4};
\node[rlab] at ($(v5)!0.500!(v19)$) {1-3};
\node[rlab] at ($(v6)!0.500!(v11)$) {3-5};
\node[rlab] at ($(v6)!0.500!(v20)$) {1-3};
\node[rlab] at ($(v7)!0.500!(v8)$) {2-3};
\node[rlab] at ($(v7)!0.500!(v12)$) {4-5};
\node[rlab] at ($(v7)!0.500!(v23)$) {1-4};
\node[rlab] at ($(v8)!0.500!(v13)$) {4-5};
\node[rlab] at ($(v8)!0.380!(v24)$) {1-4};
\node[rlab] at ($(v9)!0.500!(v10)$) {3-4};
\node[rlab] at ($(v9)!0.500!(v11)$) {2-3};
\node[rlab] at ($(v9)!0.500!(v28)$) {1-5};
\node[rlab] at ($(v10)!0.500!(v12)$) {2-4};
\node[rlab] at ($(v10)!0.500!(v29)$) {1-5};
\node[rlab] at ($(v11)!0.500!(v13)$) {3-4};
\node[rlab] at ($(v11)!0.500!(v30)$) {1-5};
\node[rlab] at ($(v12)!0.500!(v13)$) {2-3};
\node[rlab] at ($(v12)!0.500!(v31)$) {1-5};
\node[rlab] at ($(v13)!0.500!(v32)$) {1-5};
\node[rlab] at ($(v14)!0.500!(v15)$) {4-5};
\node[rlab] at ($(v14)!0.500!(v16)$) {3-4};
\node[rlab] at ($(v14)!0.500!(v21)$) {2-3};
\node[rlab] at ($(v15)!0.500!(v17)$) {3-5};
\node[rlab] at ($(v15)!0.500!(v22)$) {2-3};
\node[rlab] at ($(v16)!0.500!(v18)$) {4-5};
\node[rlab] at ($(v16)!0.500!(v25)$) {2-4};
\node[rlab] at ($(v17)!0.500!(v18)$) {3-4};
\node[rlab] at ($(v17)!0.500!(v33)$) {2-5};
\node[rlab] at ($(v18)!0.500!(v34)$) {2-5};
\node[rlab] at ($(v19)!0.500!(v20)$) {4-5};
\node[rlab] at ($(v19)!0.500!(v21)$) {1-2};
\node[rlab] at ($(v19)!0.500!(v26)$) {3-4};
\node[rlab] at ($(v20)!0.500!(v22)$) {1-2};
\node[rlab] at ($(v20)!0.500!(v35)$) {3-5};
\node[rlab] at ($(v21)!0.500!(v22)$) {4-5};
\node[rlab] at ($(v21)!0.500!(v27)$) {3-4};
\node[rlab] at ($(v22)!0.620!(v36)$) {3-5};
\node[rlab] at ($(v23)!0.500!(v24)$) {2-3};
\node[rlab] at ($(v23)!0.500!(v25)$) {1-2};
\node[rlab] at ($(v23)!0.500!(v37)$) {4-5};
\node[rlab] at ($(v24)!0.500!(v26)$) {1-3};
\node[rlab] at ($(v24)!0.500!(v38)$) {4-5};
\node[rlab] at ($(v25)!0.620!(v27)$) {2-3};
\node[rlab] at ($(v25)!0.500!(v39)$) {4-5};
\node[rlab] at ($(v26)!0.500!(v27)$) {1-2};
\node[rlab] at ($(v26)!0.500!(v40)$) {4-5};
\node[rlab] at ($(v27)!0.500!(v41)$) {4-5};
\node[rlab] at ($(v28)!0.500!(v29)$) {3-4};
\node[rlab] at ($(v28)!0.500!(v30)$) {2-3};
\node[rlab] at ($(v28)!0.500!(v33)$) {1-2};
\node[rlab] at ($(v29)!0.500!(v31)$) {2-4};
\node[rlab] at ($(v29)!0.380!(v34)$) {1-2};
\node[rlab] at ($(v30)!0.380!(v32)$) {3-4};
\node[rlab] at ($(v30)!0.500!(v35)$) {1-3};
\node[rlab] at ($(v31)!0.500!(v32)$) {2-3};
\node[rlab] at ($(v31)!0.500!(v37)$) {1-4};
\node[rlab] at ($(v32)!0.500!(v38)$) {1-4};
\node[rlab] at ($(v33)!0.500!(v34)$) {3-4};
\node[rlab] at ($(v33)!0.500!(v36)$) {2-3};
\node[rlab] at ($(v34)!0.500!(v39)$) {2-4};
\node[rlab] at ($(v35)!0.500!(v36)$) {1-2};
\node[rlab] at ($(v35)!0.500!(v40)$) {3-4};
\node[rlab] at ($(v36)!0.500!(v41)$) {3-4};
\node[rlab] at ($(v37)!0.500!(v38)$) {2-3};
\node[rlab] at ($(v37)!0.500!(v39)$) {1-2};
\node[rlab] at ($(v38)!0.500!(v40)$) {1-3};
\node[rlab] at ($(v39)!0.500!(v41)$) {2-3};
\node[rlab] at ($(v40)!0.500!(v41)$) {1-2};
\end{tikzpicture}
}
\end{center}
\caption{The graph of the $4$-dimensional associahedron: Nodes are all binary search trees on $\{1,2,3,4,5\}$. Edges connect those nodes representing trees that are a single rotation apart where the nodes of the rotation appear in the black circle on each edge. The graph is 3-colored. 
}
\end{figure}



\paragraph{Introduction.}
Associahedra are polytopes that appear in many areas of mathematics.
We refer to the survey of Pilaud, Santos, and Ziegler~\cite{MR4675114} for an overview of the developments around these polytopes in the past few decades.
We are interested in the graphs of associahedra.
%, hence from now on the word associahedron will be used to designate the graph of the associahedron.
For a given positive integer~$n$, the vertices of the $(n-1)$-dimensional associahedron are in bijection with triangulations of a convex $(n+2)$-gon, or, via planar duality, with rooted binary trees on $n$ (internal) nodes.
Hence for a given $n$, the number of vertices is $C_n = \frac{(2n)!}{(n+1)!n!}$, the $n$th Catalan number.
Edges of associahedra connect pairs of triangulations that differ by a single
\emph{diagonal-flip}, an operation that consists of 
replacing a diagonal of a quadrilateral
formed by two triangles sharing an edge by the other diagonal. Equivalently
--- and this is the representation that we will use in our proofs ---
edges connect pairs of vertices corresponding to two binary trees that differ by
a single rotation.%, the operation illustrated on Figure~\ref{f1}.

While being well-studied, the structure of the associahedron graph still remains
mysterious. It is known to have diameter $2n-6$ for sufficiently large $n$~\cite{MR4556611,MR3197650}, and
to be Hamiltonian (see for instance~\cite{MR3383157} for a generalization).
%Chromatic properties of related graphs are also investigated.
%For instance, Molnar, Pohoata, Zheng, and Zhu~\cite{MPZZ25} showed that the chromatic number of Kneser graphs of triangulations of a convex $n$-gon is $n-2$, and Barali\'c, Ivanovi\'c, and Petri\'c~\cite{MR4778792} considered coloring facets of generalized associahedra.
Chromatic properties of related graphs are also investigated~\cite{MR4778792,GargantiniPastineEtAl2026,MPZZ25}.

The natural question of the chromatic number of associahedra was posed by Fabila-Monroy, Flores-Pe\~naloza, Huemer, Hurtado, Wood, and Urrutia~\cite{DBLP:journals/dmtcs/MonroyFHHWU09} who showed a lower bound of 4 and an upper bound of $O(n/log n)$.
It is only recently that Addario-Berry, Reed, Scott, and Wood~\cite{Addario-Berry_Reed_Scott_Wood_2026} showed that the chromatic number of associahedra is at most $22500 \cdot \log_3 n + O(1) = O(\log n)$.
Their proof is elementary but quite involved. We give a simple proof that the chromatic
number of the $(n-1)$-dimensional associahedron is at most $6\log_2 n + 18$. This has since been improved to $O(\log \log n)$ \cite{oum2026doublelogarithmicupperboundchromatic} by Oum and Wood but we hope that the simple method here could give insights to further improvements.


\paragraph{Preliminaries.} We make the following definitions, which lead up to assigning any binary tree $T$ with $n$ nodes an integer color in the range 0 to $6 \lceil \log_2 n \rceil +17$.

\begin{itemize}
    \item If $x$ is a node in some tree $T$, let $s_T(x)$ be the number of nodes in the subtree of $x$ in $T$
    \item A node is \emph{left heavy} if its left subtree is larger than its right subtree; otherwise it is \emph{right heavy}.
    \item Let $L(T)$ be the set of left-heavy nodes in $T$.
    \item Let the \emph{left color} of $T$, $c_\ell(T)$ be $|L(T)| \bmod 3$.
    \item The \emph{heavy child} of a nonleaf node is its left child if it is left heavy and its right child if it is right heavy.
    \item Let $H(T)$ be the set of heavy children in $T$.
    \item Let the \emph{heavy child color} of $T$, $c_h(T)$, be $\left\lceil  \sum_{x \in H(T)} \log_2 s(x) \right\rceil \bmod \lceil \log_2n+1 \rceil$
    \item Let $P(T)$ be the set of \emph{heavy} paths, which is the graph obtained by connecting each non-leaf node to its heavy child. The \emph{heavy path depth} of a node is is its depth relative to the shallowest node (in the tree) on its heavy path. The \emph{path color}, $c_p(T)$, is the parity of the number of left heavy nodes of even path depth. 
    \item Finally, let the \emph{color} of $T$, $c(T)$ be $6c_h(T)+2c_\ell(T)+c_p(T)$. This is an integer in the range 0 to $6 \lceil \log_2 n + 1 \rceil +4+1 \leq 6 \log_2n+17$. It is designed so that if either the heavy child color or the left color or the path color differs between two trees then the color differs.
\end{itemize}

We now state the theorem:

\begin{theorem}
The chromatic number of the $(n-1)$-dimensional associahedron is at most $6\cdot \log_2 n + 18$.
\end{theorem}
\textbf{Proof.} Given the above scheme, all that remains to prove is that a single rotation always changes the color of the tree.

Consider a rotation that transforms a tree $T$ to $T'$. We assume the rotation is on two nodes $x$ and $y$ where $y$ is the right child of $x$ in $T$ and $x$ becomes the left child of $y$ in $T'$; the reverse case is symmetric. We consider 16 cases based on whether $x$ and $y$ are each left or right heavy before and after the rotation. These are illustrated in Figure~\ref{f1}. These are split into four groups, those where the left color changes, those where the heavy child color changes, and finally those cases that are impossible.

\begin{description}
    \item[\textcolor{\cac}{Cases II-V, VIII, IX, XII-XV (Left color changes)}] In these cases the number of left heavy nodes, $|L(T)|$ changes by 1 or 2, and thus $c_\ell(T)$ which is $|L(T)| \bmod 3$ changes.
    \item[\textcolor{\cbc}{Cases I and XVI (Heavy child color changes)}] For I (XVI is symmetric): Here the increase in $\sum_{x \in H} \log_2 s(x)$ is $\log_2 (|A|+|B|+1) - \log_2 |B| = \log_2 \frac{|A|+|B|+1}{|B|}$. 
    As $|A|>|B|$, $\log \frac{|A|+|B|+1}{|B|}\geq \log_2 \frac{|B|+|B|+1}{|B|}>1$. Also, $\log_2 \frac{|A|+|B|+1}{|B|}\leq \log_2 n$. Thus $\sum_{x \in H} \log_2 s(x)$ changes by at least 1 and at most $\log_2 n$, which causes $c_h(T) = \left\lceil  \sum_{x \in H} \log_2 s(x) \right\rceil \bmod \lceil \log_2n+1 \rceil$ to change.
    \item[\textcolor{\cdc}{Case X}] The key point of this case is that the heavy paths do not change, with the one exception that $x$ and $y$, which are adjacent on the same heavy path, have their positions change; $x$ remains right heavy and $y$ remains left heavy. Either $x$ before (right heavy) and $y$ after (left heavy) are at even depth on their path or $y$ before (left heavy) and $x$ after (right heavy) are at even depth; in either case the number of left heavy nodes of even path depth changes by one and thus the path color changes.
    \item[\textcolor{\ccc}{Cases VI, VII and XI (Impossible)}] These are impossible. For VI (XI is symmetric) the following contradictions occur:
    \begin{itemize}[topsep=0pt]
        \item $|A|>|B|+|C|+1>|B|$ and
        \item $|A|\leq |B|$
    \end{itemize}
    For VII:
      \begin{itemize}[topsep=0pt]
        \item $|A|>|B|+|C|+1>|C|$ and
        \item $|C|\geq |A|+|B|+1>|A|$
    \end{itemize}  
\end{description}

\paragraph{Acknowledgments} Jean Cardinal and John Iacono's research is supported by the F.R.S.–FNRS. Generative AI was used to create the Tikz figures but did not write a single word and was not used to obtain the result.




\newcommand{\ltree}[4]{%
\begin{tikzpicture}[
    every node/.style={font=\small},
    circle node/.style={circle, draw, inner sep=1pt,thick},
    triangle node/.style={
        regular polygon,
        regular polygon sides=3,
        draw=black,
        thick,
        inner sep=-0.5pt, 
%        minimum size=2mm
    },
    level distance=0.7cm,
    sibling distance=0.7cm,
    % Force all edges to go to the top of the child node
        edge from parent/.style={ultra thick},
    edge from parent path={
        (\tikzparentnode.south) -- (\tikzchildnode.north)
    }
]
%
\node[circle node] (x) {$x$}
    child {
        node[triangle node] (A) {$A$}
        edge from parent[draw=#1]
    }
    child {
        node[circle node] (y) {$y$}
        edge from parent[draw=#2]
        child {
            node[triangle node] (B) {$B$}
            edge from parent[draw=#3]
        }
        child {
            node[triangle node] (C) {$C$}
            edge from parent[draw=#4]
        }
    };
%
\end{tikzpicture}
}


\newcommand{\rtree}[4]{%
\begin{tikzpicture}[
    every node/.style={font=\small},
    circle node/.style={circle, draw, inner sep=1pt,thick},
    triangle node/.style={
        regular polygon,
        regular polygon sides=3,
        draw=black,
        inner sep=-0.5pt, 
        thick
%        minimum size=2mm
    },
    level distance=0.7cm,
    sibling distance=0.7cm,
    % Force all edges to go to the top of the child node
    edge from parent/.style={ultra thick},
    edge from parent path={
        (\tikzparentnode.south) -- (\tikzchildnode.north)
    }
]
%
\node[circle node] (x) {$y$}
    child {
        node[circle node] (y) {$x$}
        edge from parent[draw=#1]
        child {
            node[triangle node] (A) {$A$}
            edge from parent[draw=#3]
        }
        child {
            node[triangle node] (B) {$B$}
            edge from parent[draw=#4]
        }
        }
    child {
        node[triangle node] (C) {$C$}
        edge from parent[draw=#2]
    };
%
\end{tikzpicture}
}



\newcommand{\lta}{\ltree{red}{black}{red}{black}}
\newcommand{\ltb}{\ltree{red}{black}{black}{blue}}
\newcommand{\ltc}{\ltree{black}{blue}{red}{black}}
\newcommand{\ltd}{\ltree{black}{blue}{black}{blue}}

\newcommand{\rta}{\rtree{red}{black}{red}{black}}
\newcommand{\rtb}{\rtree{red}{black}{black}{blue}}
\newcommand{\rtc}{\rtree{black}{blue}{red}{black}}
\newcommand{\rtd}{\rtree{black}{blue}{black}{blue}}

\newcounter{casenum}

\newcommand{\ra}{%
{\stepcounter{casenum}%
\raisebox{42pt}{$\mathclap{\text{Case \Roman{casenum}}}$}}%
{\raisebox{20pt}{$\mathclap{\Rightarrow}$}}%
}

\newcommand{\lhc}[1]{Left heavy change: #1}



\begin{figure}
\resizebox{\textwidth}{!}{%
\begin{tabular}{|c|c|c|c|}
\hline
\lta \textcolor{\cbc}{\ra} \rta &
\lta \textcolor{\cac}{\ra}  \rtb &
\lta \textcolor{\cac}{\ra}  \rtc &
\lta \textcolor{\cac}{\ra}  \rtd 
\\
\lhc{0}
&
\lhc{-1}
&
\lhc{-1}
&
\lhc{-2}
\\ \hline
\ltb \textcolor{\cac}{\ra}  \rta &
\ltb \textcolor{\ccc}{\ra}  \rtb &
\ltb \textcolor{\ccc}{\ra}  \rtc &
\ltb \textcolor{\cac}{\ra}  \rtd 
\\
\lhc{+1}
&
Impossible
&
Impossible
&
\lhc{-1}
\\ \hline
\ltc \textcolor{\cac}{\ra}  \rta &
\ltc \textcolor{\cdc}{\ra}  \rtb &
\ltc \textcolor{\ccc}{\ra}  \rtc &
\ltc \textcolor{\cac}{\ra}  \rtd 
\\
\lhc{+1}
&
\lhc{0}
&
Impossible
&
\lhc{-1}
\\ \hline
\ltd \textcolor{\cac}{\ra} \rta &
\ltd \textcolor{\cac}{\ra}  \rtb &
\ltd \textcolor{\cac}{\ra}  \rtc &
\ltd \textcolor{\cbc}{\ra}  \rtd 
\\
\lhc{+2}
&
\lhc{+1}
&
\lhc{+1}
&
\lhc{0}
\\ \hline 
\end{tabular}
}
\caption{
The cases. We consider a left rotation on $x$ with right child $y$. A red edge indicates a left heavy node above and a blue edge indicates a right heavy node above. We consider all 16 possible cases of left and right heaviness for $x$ and $y$ before and after the rotation. 
In \textcolor{\cac}{Cases II-V, VIII, IX, XII-XV} the number of left heavy nodes change by 1 or 2 and thus the left color changes. 
In \textcolor{\cdc}{Case X} the only change in heavy paths is the relative ordering of $x$ and $y$ is swapped, each retaining their heaviness; thus the path color changes.
\textcolor{\ccc}{Cases VI, VII and XI} are impossible. 
\textcolor{\cbc}{Cases I and XVI} are possible but have no change in the number of left heavy nodes and thus require a more detailed argument and are the source of the logarithmic number of colors.
}
\label{f1}
\end{figure}


\bibliographystyle{abbrv}
\bibliography{bib}



\end{document}